\documentclass[9pt]{ctexbeamer}

% ====== 清空默认 ======
\setbeamertemplate{navigation symbols}{}
\setbeamertemplate{headline}{}
\setbeamertemplate{footline}{}

% ====== 配色 ======
\definecolor{ink}{HTML}{1B2A4A}
\definecolor{cold}{HTML}{F4F8FF}
\definecolor{ember}{HTML}{E05C2A}
\definecolor{slate}{HTML}{4B5E7D}
\definecolor{rule}{HTML}{BCC9DB}
\definecolor{accent}{HTML}{2E86DE}
\definecolor{soft}{HTML}{EBF1FA}
\definecolor{darkbg}{HTML}{0F1B33}

% ====== 顶部深色条 + 标题 ======
\setbeamercolor{section in head/foot}{fg=white,bg=ink}
\setbeamertemplate{frametitle}{
  \nointerlineskip
  \begin{beamercolorbox}[wd=\paperwidth,ht=0.85cm,dp=0.3cm,leftskip=0.5cm]{section in head/foot}
    {\large\bfseries\insertframetitle}
  \end{beamercolorbox}
}

% ====== 页脚 ======
\setbeamertemplate{footline}{
  \begin{beamercolorbox}[wd=\paperwidth,ht=0.35cm,dp=0.15cm,rightskip=0.4cm]{section in head/foot}
    {\tiny\color{rule}\insertframenumber\,/\,\inserttotalframenumber}
  \end{beamercolorbox}
}

% ====== 内容页背景 ======
\setbeamercolor{background canvas}{bg=cold}

% ====== 列表 ======
\setbeamertemplate{itemize item}{\color{ember}\rule{3pt}{3pt}}
\setbeamertemplate{itemize subitem}{\color{slate}\rule{2pt}{2pt}}

% ====== 自定义 leftblock ======
\newenvironment{leftblock}[1]{%
  \vspace{0.3cm}
  \begin{beamercolorbox}[wd=\textwidth,leftskip=0.4cm,rightskip=0.4cm]{}
    \textcolor{ember}{\rule{2.5pt}{0.55cm}}\hspace{0.2cm}
    {\bfseries\color{ink}#1}
    \par\vspace{0.1cm}
    \hspace{0.35cm}\begin{minipage}{\dimexpr\textwidth-1.2cm\relax}
}{%
    \end{minipage}
  \end{beamercolorbox}
  \vspace{0.15cm}
}

% ====== 包 ======
\usepackage{booktabs}
\usepackage{amsmath}
\usepackage{xcolor}
\usepackage{colortbl}
\usepackage{tikz}
\usepackage{pgfplots}
\pgfplotsset{compat=1.18}
\usetikzlibrary{backgrounds,calc,positioning}

% ====== pgfplots 全局 ======
\pgfplotsset{
  every axis/.append style={
    font=\tiny,
    tick label style={font=\fontsize{4.5}{5}\selectfont},
    label style={font=\tiny\color{slate}},
    grid style={dashed,gray!25,line width=0.3pt},
    axis line style={rule,line width=0.4pt},
    legend style={
      font=\fontsize{4.5}{5}\selectfont,
      legend columns=1,
      /tikz/every even column/.append style={column sep=3pt},
      draw=none,
      fill=white,fill opacity=0.7,
      inner sep=1.5pt,
      /tikz/line width=0.6pt,
    },
  },
  barplot/.style={
    ybar, width=0.93\textwidth, height=3.6cm,
    grid=major, bar width=5pt, enlarge x limits=0.12,
  },
  lineplot/.style={
    width=0.93\textwidth, height=3.2cm,
    grid=major, mark=*, mark size=1.2pt, line width=0.8pt,
  },
  narrowbar/.style={
    ybar, width=\columnwidth, height=3.0cm,
    grid=major, bar width=4pt, enlarge x limits=0.18,
  },
  narrowline/.style={
    width=\columnwidth, height=2.8cm,
    grid=major, mark=*, mark size=1pt, line width=0.7pt,
  },
}

% ====== 快捷命令 ======
\newcommand{\tb}{\textcolor{ember}{\textbf{【待补充】}}}
\newcommand{\hi}[1]{\textcolor{ember}{\textbf{#1}}}
\newcommand{\blu}[1]{\textcolor{accent}{\textbf{#1}}}

\begin{document}

% ====== Slide 1: 封面 ======
{
\setbeamercolor{background canvas}{bg=darkbg}
\setbeamertemplate{footline}{}
\setbeamertemplate{frametitle}{}

\begin{frame}[plain]

\begin{tikzpicture}[remember picture,overlay]
  \shade[top color=darkbg,bottom color=ink] (current page.south west) rectangle (current page.north east);
  \foreach \y in {0,2,...,16} {
    \foreach \x in {0,2,...,12} {
      \pgfmathsetmacro\r{random()}
      \pgfmathsetmacro\c{int(\r < 0.42 ? 1 : 0)}
      \node[text=white,opacity=0.09,font=\ttfamily\tiny] at ($(current page.north west)+(1.2+0.85*\x,-0.9-0.6*\y)$) {\c};
    }
  }
  \draw[ember,line width=0.6pt,opacity=0.5] (current page.north west) -- ($(current page.north west)!0.4!(current page.south east)$);
  \draw[accent,line width=0.4pt,opacity=0.35] ($(current page.north east)!0.15!(current page.south east)$) -- ($(current page.south east)!0.3!(current page.south west)$);
\end{tikzpicture}

\begin{center}
\vspace*{1.2cm}
{\fontsize{24}{30}\selectfont\color{white}\textbf{多种算法求解\\[0.25cm]0/1背包问题}}

\vspace{0.25cm}
\textcolor{ember}{\rule{2.8cm}{1.6pt}}

\vspace{0.6cm}
{\large\color{cold}性能比较实验}

\vspace{1.2cm}
{\normalsize\color{rule}贾静潇 \quad 张诗淇}

\vspace{0.25cm}
{\footnotesize\color{rule}算法分析与设计 · 2.3 问题求解类}

\vspace{0.4cm}
{\tiny\color{rule!60}数据集: FSU · Pi · JJ \quad 算法: 贪心 · DP · 回溯 · 分支限界}
\end{center}
\end{frame}
}

% ====== Slide 2: 问题定义 ======
\begin{frame}{0/1背包问题}

\begin{leftblock}{形式化定义}
给定 $n$ 个物品（重量 $w_i$、价值 $v_i$）和容量 $W$，
\[\text{max}\sum_{i\in S} v_i \qquad \text{s.t.}\;\sum_{i\in S} w_i \leq W,\; x_i \in \{0,1\}\]
\end{leftblock}

\vspace{0.2cm}

\begin{columns}[T]
\begin{column}{0.5\textwidth}
{\small\color{ink}\textbf{理论性质}}\\[0.1cm]
{\footnotesize\color{slate}NP-hard · 最优子结构 · 伪多项式可解}
\end{column}
\begin{column}{0.5\textwidth}
{\small\color{ink}\textbf{应用场景}}\\[0.1cm]
{\footnotesize\color{slate}资源分配 · 投资决策 · 货物装载 · 带宽分配}
\end{column}
\end{columns}

\vspace{0.3cm}
\begin{itemize}
  \item 每件物品\hi{要么整取、要么不取}，不可分割
  \item 物品间\hi{无先后顺序}约束，仅受容量限制
  \item 目标：在 $2^n$ 种组合中找到满足容量约束的最大价值子集
\end{itemize}

\end{frame}

% ====== Slide 3: 四种算法总览 ======
\begin{frame}{四种求解策略}

\vspace{-0.1cm}
\begin{table}
\renewcommand{\arraystretch}{1.25}
\small
\begin{tabular}{p{1.8cm}p{2.5cm}p{2.2cm}p{2.2cm}}
\toprule
\textbf{算法} & \textbf{核心思路} & \textbf{时间} & \textbf{最优？} \\
\midrule
\hi{贪心}      & 性价比排序，依次塞入   & $O(n\log n)$  & \hi{否} \\
\blu{动态规划} & 自底向上递推填表       & $O(nW)$       & 是 \\
回溯          & 深搜 + 上界剪枝       & 最坏 $O(2^n)$ & 是 \\
分支限界      & 优先队列 + 上界剪枝   & 最坏 $O(2^n)$ & 是 \\
\bottomrule
\end{tabular}
\end{table}

\vspace{0.15cm}
\begin{center}
\begin{tikzpicture}[font=\footnotesize]
  \node[draw=rule,fill=soft,text width=8cm,align=center,rounded corners=3pt,inner sep=6pt] {
    {\color{ink}\textbf{同一问题 · 四种策略 · 不同瓶颈}}\\[0.08cm]
    {\color{slate}\small 贪心怕数据分布 \quad DP 怕大容量 \quad 搜索怕大物品数}
  };
\end{tikzpicture}
\end{center}

\end{frame}

% ====== Slide 4: 贪心 ======
\begin{frame}{贪心算法}

\begin{columns}[T]
\begin{column}{0.5\textwidth}

\begin{leftblock}{策略}
$r_i \gets v_i/w_i$，按 $r_i$ \hi{降序}排列，依次尝试放入——够就塞，跳过不回头。
\end{leftblock}

\vspace{0.1cm}
\begin{itemize}
  \item 时间 $O(n\log n)$ · 空间 $O(n)$
  \item 与 $W$ \hi{完全无关}
  \item 大 $n$ 下依然毫秒级
\end{itemize}

\end{column}
\begin{column}{0.5\textwidth}

\begin{leftblock}{失效示例}
容量 10
\vspace{0.05cm}
\begin{tabular}{cccc}
\hi{A} & w=6 & v=9  & $r$=1.5 \\
B     & w=5 & v=7  & $r$=1.4 \\
C     & w=5 & v=7  & $r$=1.4 \\
\end{tabular}
\par\vspace{0.1cm}
贪心拿 A：价值 9（剩余 4）\\
\blu{最优 B+C：价值 14}
\end{leftblock}

\end{column}
\end{columns}

\vspace{0.2cm}
\begin{center}
{\small\color{slate}每次选当前最优 → 局部最优不保证全局最优 → \hi{需要后悔，但不后悔}}
\end{center}

\end{frame}

% ====== Slide 5: DP ======
\begin{frame}{动态规划}

\begin{columns}[T]
\begin{column}{0.5\textwidth}
\begin{leftblock}{状态与转移}
\[\boxed{dp[w] = \max\{\, dp[w],\; dp[w-w_i] + v_i \,\}}\]
$dp[w]$ = 容量 $w$ 时的最大价值
\end{leftblock}

\vspace{0.1cm}
\begin{itemize}
  \item 内层 $w$ 必须\hi{倒序}
  \item 正序 → 物品可重复 → 完全背包
\end{itemize}
\end{column}

\begin{column}{0.5\textwidth}
\begin{leftblock}{复杂度}
\textbf{时间：} \hi{$O(nW)$} 伪多项式 \\
\textbf{空间：} $O(W)$（一维优化）

\vspace{0.12cm}
{\footnotesize\color{slate}
$n$=2000, $W\approx250$万：数十亿次状态转移 → 约10s \\
$W$继续放大时：dp数组线性膨胀 → 内存压力明显
}
\end{leftblock}
\end{column}
\end{columns}

\vspace{0.2cm}
\begin{center}
{\small\color{slate}以 $W$ 为代价换最优解 · 当 $W$ 失控时从优势变为劣势}
\end{center}

\end{frame}

% ====== Slide 6: 回溯与分支限界 ======
\begin{frame}{回溯与分支限界}

\begin{columns}[T]
\begin{column}{0.5\textwidth}
\begin{leftblock}{回溯 Backtracking}
深度优先遍历“选/不选”解空间树

\vspace{0.08cm}
{\footnotesize\color{slate}
可行性剪枝：超重立即返回 \\
上界剪枝：当前价值 + 剩余上界 $\leq$ 最优值则停止
}
\end{leftblock}

\vspace{0.05cm}
{\footnotesize\color{slate}
\textbf{复杂度：}最坏 $O(2^n)$，空间 $O(n)$。\\
\textbf{适合：}小规模精确求解、验证最优值。
}
\end{column}
\begin{column}{0.5\textwidth}
\begin{leftblock}{分支限界 Branch \\ Bound}
优先队列维护活节点，优先扩展上界最高的节点

\vspace{0.08cm}
{\footnotesize\color{slate}
用分数背包思想估计上界 \\
上界不足的节点直接丢弃
}
\end{leftblock}

\vspace{0.05cm}
{\footnotesize\color{slate}
\textbf{复杂度：}最坏 $O(2^n)$，空间最坏 $O(2^n)$。\\
\textbf{适合：}剪枝有效的小/中规模精确搜索。
}
\end{column}
\end{columns}

\vspace{0.25cm}
\begin{center}
\begin{tikzpicture}[font=\footnotesize]
  \node[draw=rule,fill=soft,text width=8.3cm,align=center,rounded corners=3pt,inner sep=5pt] {
    {\color{ink}\textbf{实验处理：}}FSU 与 Pi-S 完整运行；Pi-M/Pi-L 与 JJ 因 $n$ 过大记录为 \hi{SKIP}。\\[0.04cm]
    {\color{slate}\scriptsize 这不是代码缺失，而是指数级搜索在大规模输入下的实际限制。}
  };
\end{tikzpicture}
\end{center}

\end{frame}

% ====== Slide 7: 数据集 ======
\begin{frame}{三个数据源}

\vspace{-0.2cm}
\begin{table}
\renewcommand{\arraystretch}{1.2}
\footnotesize
\begin{tabular}{p{1.5cm}p{2.3cm}p{3.8cm}r}
\toprule
\textbf{数据源} & \textbf{来源} & \textbf{特征} & \textbf{数量} \\
\midrule
\hi{FSU}  & Florida State U. & 官方基准，$n$=5--24，附最优解对照 & 8 \\
\blu{Pi}  & 自生成(Pisinger) & 3种相关性 $\times$ 3种规模，固定种子 & 9 \\
\hi{JJ}  & JorikJooken (GH) & 学术标准集，$n$=400--1000，$W$=100万 & 6 \\
\bottomrule
\end{tabular}
\end{table}

\vspace{0.1cm}
\begin{center}
\begin{tikzpicture}[font=\footnotesize]
  \node[draw=accent,fill=accent!8,text width=2.8cm,align=center,rounded corners=3pt,inner sep=4pt] (f) {\textbf{FSU}\\[0.03cm]\tiny 正确性验证};
  \node[draw=accent,fill=accent!8,text width=2.8cm,align=center,rounded corners=3pt,inner sep=4pt,right=0.3cm of f] (p) {\textbf{Pi}\\[0.03cm]\tiny 控制变量分析};
  \node[draw=accent,fill=accent!8,text width=2.8cm,align=center,rounded corners=3pt,inner sep=4pt,right=0.3cm of p] (j) {\textbf{JJ}\\[0.03cm]\tiny 大规模实战测试};
  \draw[->,ember,line width=1pt] (f) -- (p);
  \draw[->,ember,line width=1pt] (p) -- (j);
\end{tikzpicture}
\end{center}

\end{frame}


% ====== Slide 8: FSU 四算法同表对比 ======
\begin{frame}{FSU：四种算法同表比较}

\vspace{-0.18cm}
\begin{table}
\renewcommand{\arraystretch}{1.28}
\footnotesize
\begin{tabular}{p{1.15cm}r rrrr r}
\toprule
\textbf{实例} & \textbf{贪心} & \textbf{DP} & \textbf{回溯} & \textbf{分支限界} & \textbf{最优值} & \textbf{贪心Gap} \\
\midrule
P01 & 309 & 309 & 309 & 309 & 309 & 0.00\% \\
P06 & \hi{1478} & \blu{1735} & \blu{1735} & \blu{1735} & 1735 & \hi{14.81\%} \\
P08 & \hi{13415886} & \blu{13549094} & \blu{13549094} & \blu{13549094} & 13549094 & 0.98\% \\
\bottomrule
\end{tabular}
\end{table}

\vspace{0.15cm}
\begin{columns}[T]
\begin{column}{0.48\textwidth}
\begin{leftblock}{同一张表看四种算法}
{\footnotesize\color{slate}
FSU 数据规模小，四种算法都可运行。DP、回溯、分支限界全部命中最优，贪心在 P06 出现明显偏差。}
\end{leftblock}
\end{column}
\begin{column}{0.48\textwidth}
\begin{leftblock}{分析重点}
{\footnotesize\color{slate}
P06 说明“性价比排序”不等于全局最优；P08 则体现容量 $W$ 对 DP 的影响。}
\end{leftblock}
\end{column}
\end{columns}

\vspace{0.1cm}
\begin{center}
{\footnotesize\color{ink}\textbf{结论：}小规模场景下，搜索类算法能作为 DP 的精确验证；贪心只能作为快速基线。}
\end{center}

\end{frame}

% ====== Slide 9: FSU 四算法耗时 ======
\begin{frame}{FSU：四种算法耗时比较}

\vspace{-0.28cm}
\begin{center}
{\small\color{ink}\textbf{代表实例耗时对比（log 坐标，ms）}}
\begin{tikzpicture}
\begin{axis}[
  ybar, width=0.94\textwidth, height=4.55cm,
  ymode=log, log basis y={10},
  grid=major, bar width=6pt, enlarge x limits=0.22,
  ylabel={耗时(ms)},
  symbolic x coords={P06,P08},
  xtick=data,
  ymin=0.01,
  legend style={font=\fontsize{5}{5.5}\selectfont,at={(0.5,-0.22)},anchor=north,legend columns=4,draw=none,fill=none},
  x tick label style={font=\fontsize{6}{6.5}\selectfont},
]
  \addplot[fill=accent!75,draw=accent] coordinates {(P06,0.069)(P08,757.697)};
  \addlegendentry{DP}
  \addplot[fill=ember!75,draw=ember] coordinates {(P06,0.161)(P08,0.081)};
  \addlegendentry{贪心}
  \addplot[fill=slate!65,draw=slate] coordinates {(P06,0.070)(P08,0.306)};
  \addlegendentry{回溯}
  \addplot[fill=ink!70,draw=ink] coordinates {(P06,0.271)(P08,1.378)};
  \addlegendentry{分支限界}
\end{axis}
\end{tikzpicture}
\end{center}

\vspace{0.08cm}
\begin{itemize}
  \item P08 的 $n=24$，但 $W=6,404,180$，DP 需要枚举容量状态，因此耗时显著增加。
  \item 回溯/分支限界主要受 $n$ 和剪枝效果影响，\hi{不直接按 $W$ 扩张}，所以在 P08 仍可运行。
\end{itemize}

\end{frame}

% ====== Slide 10: Pi-S 四算法比较 ======
\begin{frame}{Pi-S：四种算法完整比较}

\vspace{-0.28cm}
\begin{center}
{\small\color{ink}\textbf{S规模三类数据耗时对比（log 坐标，ms）}}
\begin{tikzpicture}
\begin{axis}[
  ybar, width=0.94\textwidth, height=4.55cm,
  ymode=log, log basis y={10},
  grid=major, bar width=5pt, enlarge x limits=0.17,
  ylabel={耗时(ms)},
  symbolic x coords={不相关S,弱相关S,强相关S},
  xtick=data,
  ymin=0.05,
  legend style={font=\fontsize{5}{5.5}\selectfont,at={(0.5,-0.22)},anchor=north,legend columns=4,draw=none,fill=none},
  x tick label style={font=\fontsize{5.5}{6}\selectfont},
]
  \addplot[fill=accent!75,draw=accent] coordinates {(不相关S,16.967)(弱相关S,3.882)(强相关S,1.558)};
  \addlegendentry{DP}
  \addplot[fill=ember!75,draw=ember] coordinates {(不相关S,23.471)(弱相关S,0.173)(强相关S,0.083)};
  \addlegendentry{贪心}
  \addplot[fill=slate!65,draw=slate] coordinates {(不相关S,23.391)(弱相关S,0.689)(强相关S,0.218)};
  \addlegendentry{回溯}
  \addplot[fill=ink!70,draw=ink] coordinates {(不相关S,34.255)(弱相关S,1.092)(强相关S,1.549)};
  \addlegendentry{分支限界}
\end{axis}
\end{tikzpicture}
\end{center}

\vspace{0.05cm}
\begin{columns}[T]
\begin{column}{0.52\textwidth}
{\footnotesize\color{slate}S规模 $n=50$：四种算法全部 \blu{OK}，DP、回溯、分支限界均命中最优。}
\end{column}
\begin{column}{0.45\textwidth}
{\footnotesize\color{slate}M/L规模 $n=500/2000$：搜索空间过大，回溯和分支限界记录为 \hi{SKIP}。}
\end{column}
\end{columns}

\end{frame}

% ====== Slide 11: 搜索类算法状态矩阵 ======
\begin{frame}{搜索类算法：从可运行到不可运行}

\vspace{-0.15cm}
\begin{center}
\renewcommand{\arraystretch}{1.35}
\small
\begin{tabular}{p{2.2cm}p{1.5cm}p{1.8cm}p{1.8cm}p{2.35cm}}
\toprule
\textbf{数据组} & \textbf{$n$} & \textbf{回溯} & \textbf{分支限界} & \textbf{说明} \\
\midrule
FSU & 5--24 & \blu{OK} & \blu{OK} & 全部命中最优 \\
Pi-S & 50 & \blu{OK} & \blu{OK} & 小规模完整比较 \\
Pi-M/L & 500/2000 & \hi{SKIP} & \hi{SKIP} & 指数搜索空间过大 \\
JJ & 400/1000 & \hi{SKIP} & \hi{SKIP} & 大规模实例不可完整搜索 \\
\bottomrule
\end{tabular}
\end{center}

\vspace{0.25cm}
\begin{columns}[T]
\begin{column}{0.48\textwidth}
\begin{leftblock}{回溯}
{\footnotesize\color{slate}递归深度约 $O(n)$，内存压力较低；但节点数最坏 $2^n$，规模扩大后迅速失控。}
\end{leftblock}
\end{column}
\begin{column}{0.48\textwidth}
\begin{leftblock}{分支限界}
{\footnotesize\color{slate}优先队列让搜索更有方向；但活节点过多时，空间最坏也可能接近指数级。}
\end{leftblock}
\end{column}
\end{columns}

\vspace{0.08cm}
\begin{center}
{\small\color{ink}\textbf{结论：搜索类算法保证最优，但“保证最优”不等于“大规模可运行”。}}
\end{center}

\end{frame}

% ====== Slide 12: Pi 结果 ======
\begin{frame}{Pi：规模扩展与贪心偏差}

\vspace{-0.28cm}
\begin{columns}[T]
\begin{column}{0.49\textwidth}
{\small\color{ink}\textbf{DP 耗时趋势（log）}}\\[-0.03cm]
\begin{tikzpicture}
\begin{axis}[
  width=\columnwidth, height=3.25cm,
  ymode=log, log basis y={10},
  grid=major, mark=*, mark size=1pt, line width=0.7pt,
  ylabel={耗时(ms)},
  xtick={0,1,2}, xticklabels={S,M,L},
  legend style={font=\fontsize{4.3}{4.8}\selectfont,at={(0.5,-0.25)},anchor=north,legend columns=3,draw=none,fill=none},
]
  \addplot[accent] coordinates {(0,16.967)(1,528.519)(2,10056.173)}; \addlegendentry{不相关}
  \addplot[ember] coordinates {(0,3.882)(1,657.953)(2,10663.162)}; \addlegendentry{弱相关}
  \addplot[ink] coordinates {(0,1.558)(1,739.220)(2,11416.532)}; \addlegendentry{强相关}
\end{axis}
\end{tikzpicture}
\end{column}
\begin{column}{0.49\textwidth}
{\small\color{ink}\textbf{贪心 Gap 按模式}}\\[-0.03cm]
\begin{tikzpicture}
\begin{axis}[
  ybar, width=\columnwidth, height=3.25cm,
  grid=major, bar width=4pt, enlarge x limits=0.18,
  ylabel={Gap(\%)},
  symbolic x coords={不相关,弱相关,强相关},
  xtick=data,
  legend style={font=\fontsize{4.3}{4.8}\selectfont,at={(0.5,-0.25)},anchor=north,legend columns=3,draw=none,fill=none},
  x tick label style={font=\fontsize{4.3}{4.8}\selectfont},
]
  \addplot[fill=accent!55,draw=accent] coordinates {(不相关,0.00)(弱相关,0.00)(强相关,0.85)}; \addlegendentry{S}
  \addplot[fill=accent!80,draw=accent] coordinates {(不相关,0.01)(弱相关,0.04)(强相关,0.12)}; \addlegendentry{M}
  \addplot[fill=accent,draw=accent] coordinates {(不相关,0.00)(弱相关,0.00)(强相关,0.02)}; \addlegendentry{L}
\end{axis}
\end{tikzpicture}
\end{column}
\end{columns}

\vspace{0.18cm}
{\footnotesize\color{slate}DP 从 S 到 L 耗时显著上升；贪心在 Pi 中整体接近最优，但强相关数据中排序区分度降低，gap 更容易出现。}

\end{frame}

% ====== Slide 13: JJ 结果 ======
\begin{frame}{JJ：大规模实例下的现实选择}

\vspace{-0.28cm}
\begin{center}
{\small\color{ink}\textbf{DP vs 贪心耗时（log 坐标，ms）}}
\begin{tikzpicture}
\begin{axis}[
  ybar, width=0.94\textwidth, height=4.45cm,
  ymode=log, log basis y={10},
  grid=major, bar width=6pt, enlarge x limits=0.11,
  ylabel={耗时(ms)},
  symbolic x coords={400-0.1,400-0.2,400-0.3,1000-0.1,1000-0.2,1000-0.3},
  xtick=data,
  ymin=0.1,
  legend style={font=\fontsize{5}{5.5}\selectfont,at={(0.5,-0.20)},anchor=north,legend columns=2,draw=none,fill=none},
  x tick label style={font=\fontsize{4.8}{5.2}\selectfont},
]
  \addplot[fill=ember!80,draw=ember] coordinates {(400-0.1,27.783)(400-0.2,0.516)(400-0.3,0.990)(1000-0.1,2.051)(1000-0.2,1.860)(1000-0.3,0.959)};
  \addlegendentry{贪心}
  \addplot[fill=accent!80,draw=accent] coordinates {(400-0.1,788.311)(400-0.2,750.138)(400-0.3,727.072)(1000-0.1,1556.003)(1000-0.2,1616.109)(1000-0.3,1632.430)};
  \addlegendentry{DP}
\end{axis}
\end{tikzpicture}
\end{center}

\vspace{0.02cm}
\begin{itemize}
  \item DP 6 个 JJ 实例全部命中最优；$n=1000$ 时耗时约 \hi{1.6s}。
  \item 贪心 gap 最大仅 \hi{0.04\%}，但仍不具备理论最优保证。
  \item 回溯/分支限界面对 $n=400/1000$ 的指数搜索树，默认记录为 \hi{SKIP}。
\end{itemize}

\end{frame}
% ====== Slide 11: 核心洞察 ======
\begin{frame}{核心洞察：$n$ 与 $W$ 的博弈}

\vspace{-0.15cm}
\begin{center}
\renewcommand{\arraystretch}{1.8}
\large
\begin{tabular}{c|c|c|}
\multicolumn{1}{c}{} & \multicolumn{1}{c}{\textbf{$n$ 小}} & \multicolumn{1}{c}{\textbf{$n$ 大}} \\
\cline{2-3}
\textbf{$W$ 小} & \color{slate}四类算法均可比较 & \blu{DP/贪心更现实} \\
\cline{2-3}
\textbf{$W$ 大} & \hi{搜索类可避开 $W$ 维度} & \hi{贪心最快，DP 受 $W$ 限制} \\
\cline{2-3}
\end{tabular}
\end{center}

\vspace{0.25cm}

\begin{columns}[T]
\begin{column}{0.52\textwidth}
{\footnotesize\color{slate}
DP 的敌人是 \hi{$W$}：容量越大，状态数组越大。\\
回溯的敌人是 \hi{$n$}：每增加一个物品，理论搜索空间翻倍。\\
贪心两者都不太怕，但以\hi{最优性不保证}为代价。
}
\end{column}
\begin{column}{0.44\textwidth}
\begin{center}
\textcolor{ember}{\rule{2cm}{0.5pt}}\\[0.2cm]
{\large\color{ink}\textbf{没有万能算法\\只有合适选型}}\\[0.2cm]
\textcolor{ember}{\rule{2cm}{0.5pt}}
\end{center}
\end{column}
\end{columns}

\end{frame}

% ====== Slide 12: 复杂度总表 ======
\begin{frame}{四种算法复杂度对比}

\vspace{-0.2cm}
\begin{table}
\renewcommand{\arraystretch}{1.3}
\small
\begin{tabular}{p{1.9cm}p{1.8cm}p{1.8cm}p{1.8cm}p{1.8cm}}
\toprule
\textbf{维度} & \textbf{贪心} & \textbf{DP} & \textbf{回溯} & \textbf{分支限界} \\
\midrule
时间复杂度 & $O(n\log n)$ & \hi{$O(nW)$} & $O(2^n)$ & $O(2^n)$ \\
空间复杂度 & $O(n)$ & \hi{$O(W)$} & $O(n)$ & 最坏 $O(2^n)$ \\
保证最优？ & \hi{否} & 是 & 是 & 是 \\
敏感于 & 数据分布 & $W$ & $n$ & $n$ / 上界 \\
大规模可行？ & \blu{强} & 取决于 $W$ & \hi{弱} & \hi{弱} \\
\bottomrule
\end{tabular}
\end{table}

\vspace{0.15cm}
\begin{center}
{\large\color{ink}贪心：快但不准 \quad$\mid$\quad DP：准但怕大W \quad$\mid$\quad 搜索类：准但怕大n}\\[0.2cm]
{\footnotesize\color{slate}四种算法构成连续谱系——从极速近似到精确最优，选型取决于 $n$、$W$ 和精度约束。}
\end{center}

\end{frame}

% ====== Slide 13: 总结 ======
\begin{frame}{总结}

\vspace{-0.1cm}
\begin{enumerate}
  \item 0/1背包问题以\hi{简洁的结构承载了丰富的算法策略差异}

  \item \hi{贪心} — 毫秒级 · 与 $W$ 无关 · gap 最高 15\% · \blu{实时场景首选}

  \item \blu{DP} — 始终最优 · $O(nW)$ 伪多项式 · 大容量下开销明显 · \blu{精度优先}

  \item 回溯/分支限界 — 小规模精确可行 · 大规模受 $2^n$ 限制 · FSU/Pi-S 全部最优

  \item \textbf{选型公式} — 看 $n$、看 $W$、看精度需求 → \hi{速度、空间、最优性三者取舍}
\end{enumerate}

\vspace{0.3cm}
\begin{center}
{\large\color{ink}\textbf{时间与最优性的折中，是算法设计的永恒主题。}}
\end{center}

\vspace{0.15cm}
\begin{flushright}
{\footnotesize\color{rule}谢谢！}
\end{flushright}

\end{frame}

\end{document}
